
\documentclass[a4paper,11pt]{article}
\usepackage[a4paper, margin=8em]{geometry}

% usa i pacchetti per la scrittura in italiano
\usepackage[french,italian]{babel}
\usepackage[T1]{fontenc}
\usepackage[utf8]{inputenc}
\frenchspacing 

% usa i pacchetti per la formattazione matematica
\usepackage{amsmath, amssymb, amsthm, amsfonts}

% usa altri pacchetti
\usepackage{gensymb}
\usepackage{hyperref}
\usepackage{standalone}

% imposta il titolo
\title{Appunti Elettronica Digitale}
\author{Luca Seggiani}
\date{2026}

% imposta lo stile
% usa helvetica
\usepackage[scaled]{helvet}
% usa palatino
\usepackage{palatino}
% usa un font monospazio guardabile
\usepackage{lmodern}

\renewcommand{\rmdefault}{ppl}
\renewcommand{\sfdefault}{phv}
\renewcommand{\ttdefault}{lmtt}

% disponi il titolo
\makeatletter
\renewcommand{\maketitle} {
	\begin{center} 
		\begin{minipage}[t]{.8\textwidth}
			\textsf{\huge\bfseries \@title} 
		\end{minipage}%
		\begin{minipage}[t]{.2\textwidth}
			\raggedleft \vspace{-1.65em}
			\textsf{\small \@author} \vfill
			\textsf{\small \@date}
		\end{minipage}
		\par
	\end{center}

	\thispagestyle{empty}
	\pagestyle{fancy}
}
\makeatother

% disponi teoremi
\usepackage{tcolorbox}
\newtcolorbox[auto counter, number within=section]{theorem}[2][]{%
	colback=blue!10, 
	colframe=blue!40!black, 
	sharp corners=northwest,
	fonttitle=\sffamily\bfseries, 
	title=~\thetcbcounter: #2, 
	#1
}

% disponi definizioni
\newtcolorbox[auto counter, number within=section]{definition}[2][]{%
	colback=red!10,
	colframe=red!40!black,
	sharp corners=northwest,
	fonttitle=\sffamily\bfseries,
	title=~\thetcbcounter: #2,
	#1
}

% U.D.M
\newcommand{\amp}{\ensuremath{\mathrm{A}}}
\newcommand{\volt}{\ensuremath{\mathrm{V}}}
\newcommand{\meter}{\ensuremath{\mathrm{m}}}
\newcommand{\second}{\ensuremath{\mathrm{s}}}
\newcommand{\farad}{\ensuremath{\mathrm{F}}}
\newcommand{\henry}{\ensuremath{\mathrm{H}}}
\newcommand{\siemens}{\ensuremath{\mathrm{S}}}

% circuiti
\usepackage{circuitikz}
\usetikzlibrary{babel}

% disegni
\usepackage{pgfplots}
\pgfplotsset{width=10cm,compat=1.9}

% disponi codice
\usepackage{listings}
\usepackage[table]{xcolor}

\definecolor{codegreen}{rgb}{0,0.6,0}
\definecolor{codegray}{rgb}{0.5,0.5,0.5}
\definecolor{codepurple}{rgb}{0.58,0,0.82}
\definecolor{backcolour}{rgb}{0.95,0.95,0.92}

\lstdefinestyle{codestyle}{
		backgroundcolor=\color{black!5}, 
		commentstyle=\color{codegreen},
		keywordstyle=\bfseries\color{magenta},
		numberstyle=\sffamily\tiny\color{black!60},
		stringstyle=\color{green!50!black},
		basicstyle=\ttfamily\footnotesize,
		breakatwhitespace=false,         
		breaklines=true,                 
		captionpos=b,                    
		keepspaces=true,                 
		numbers=left,                    
		numbersep=5pt,                  
		showspaces=false,                
		showstringspaces=false,
		showtabs=false,                  
		tabsize=2
}

\lstdefinestyle{shellstyle}{
		backgroundcolor=\color{black!5}, 
		basicstyle=\ttfamily\footnotesize\color{black}, 
		commentstyle=\color{black}, 
		keywordstyle=\color{black},
		numberstyle=\color{black!5},
		stringstyle=\color{black}, 
		showspaces=false,
		showstringspaces=false, 
		showtabs=false, 
		tabsize=2, 
		numbers=none, 
		breaklines=true
}

\lstdefinelanguage{javascript}{
	keywords={typeof, new, true, false, catch, function, return, null, catch, switch, var, if, in, while, do, else, case, break},
	keywordstyle=\color{blue}\bfseries,
	ndkeywords={class, export, boolean, throw, implements, import, this},
	ndkeywordstyle=\color{darkgray}\bfseries,
	identifierstyle=\color{black},
	sensitive=false,
	comment=[l]{//},
	morecomment=[s]{/*}{*/},
	commentstyle=\color{purple}\ttfamily,
	stringstyle=\color{red}\ttfamily,
	morestring=[b]',
	morestring=[b]"
}

% disponi sezioni
\usepackage{titlesec}

\titleformat{\section}
	{\sffamily\Large\bfseries} 
	{\thesection}{1em}{} 
\titleformat{\subsection}
	{\sffamily\large\bfseries}   
	{\thesubsection}{1em}{} 
\titleformat{\subsubsection}
	{\sffamily\normalsize\bfseries} 
	{\thesubsubsection}{1em}{}

% disponi alberi
\usepackage{forest}

\forestset{
	rectstyle/.style={
		for tree={rectangle,draw,font=\large\sffamily}
	},
	roundstyle/.style={
		for tree={circle,draw,font=\large}
	}
}

% disponi algoritmi
\usepackage{algorithm}
\usepackage{algorithmic}
\makeatletter
\renewcommand{\ALG@name}{Algoritmo}
\makeatother

% disponi numeri di pagina
\usepackage{fancyhdr}
\fancyhf{} 
\fancyfoot[L]{\sffamily{\thepage}}

\makeatletter
\fancyhead[L]{\raisebox{1ex}[0pt][0pt]{\sffamily{\@title \ \@date}}} 
\fancyhead[R]{\raisebox{1ex}[0pt][0pt]{\sffamily{\@author}}}
\makeatother

\begin{document}
% sezione (data)
\section{Lezione del 05-03-26}

% stili pagina
\thispagestyle{empty}
\pagestyle{fancy}

% testo
\subsection{Richiami di Elettrotecnica}
Dedichiamo questa lezione ad alcuni richiami di conoscenze passate. Notiamo che appunti più dettagliati di elettrotecnica possono trovarsi a \url{https://raw.githubusercontent.com/seggiani-luca/appunti-et/e63e7e9ac9d57bd7b3f690140c8b460d91e0b08e/master/master.pdf}. Qui faremo solo un breve ripasso degli argomenti più importanti, cioè:
\begin{itemize}
	\item Leggi di \textbf{Ohm} (già viste in 2.1);
	\item Leggi di \textbf{Kirchoff} sull'analisi circuitale;
	\item Teoremi di \textbf{Thevenin} e \textbf{Norton};
	\item Componenti circuitali passivi, cioè \textit{resistenze}, \textit{capacità} ed \textit{induttanze}.
\end{itemize}

\par\smallskip

\noindent
\textbf{\textsf{Leggi di Ohm}}

\noindent
In particolare, in 2.1 abbiamo ritrovato le leggi di Ohm, che riportiamo qui per completezza:
$$
v = Ri, \ i = G v \quad \textit{($1^a$ legge di Ohm)}
$$
$$
R = \frac{1}{G} = \rho \frac{l}{A}, \quad \textit{($2^a$ legge di Ohm)}
$$
Queste stabiliscono che la tensione $v$ ai capi di un conduttore è proporzionale alla corrente $i$ che percorre tale conduttore, di una costante $R$, detta \textit{resistenza} (o il suo inverso $G$, detto \textit{conduttanza}).

\par\smallskip

\noindent
\textbf{\textsf{Leggi di Kirchoff}}

\noindent
Vediamo quindi le leggi di analisi circuitale, a partire da:
\begin{theorem}{Prima legge di Kirchoff (ai nodi)}
La somma algebrica (con segno) delle correnti entranti in un qualsiasi nodo di un circuito è uguale a zero, ergo:
$$
\sum_{n=1}^N i_n = 0
$$
\end{theorem}

Per questa come sempre ricordiamo che le correnti hanno segno, per cui è importante prima stabilire i versi di tutte le correnti e quindi applicare la legge per evitare di incorrere in errori di segno.

\newpage

Vediamo quindi la seconda legge:
\begin{theorem}{Seconda legge di Kirchoff (alle maglie)}
La somma algebrica delle cadute di tensione lungo un qualunque percorso chiuso di un circuito (detto \textit{maglia}) è uguale a zero, ergo:
$$
\sum_{n=1}^N v_n = 0
$$
\end{theorem}

Spesso, in circuiti su cui sono presenti generatori, questo significa che la somma dei generatori di tensione è uguale alla caduta di potenziale sui terminali degli altri componenti. In altre parole:
	 "La fem dei generatori è uguale alla somma algebrica delle cadute di tensione negli elementi circuitali passivi", ergo:
 $$
	\sum_{i=0}^N E_i = \sum_{j=0}^M i R_i
 $$
 per generatori $E_i$ e resistenze $R_I$ su una maglia percorsa da corrente $i$.

 \par\smallskip

\noindent
\textbf{\textsf{Partitore di tensione}}

\noindent
Abbiamo che la resistenza equivalente di resistenze serie è banalmente:
$$
R_{eq} = R_1 + R_2 + ... + R_j + ... + R_n
$$

Possiamo usare questa legge per studiare il \textit{partitore di tensione}:
\begin{center}
\begin{circuitikz} 
	\draw
  (0,2) to[V, v=$v$] (0,12)
  -- (2,12) 
  to[R, l=$R_1$] (2,10) 
  to[R, l=$R_2$] (2,8);

	\draw
	(2,7) to[R, l=$R_j$] (2,5);

	\draw
	(2,4) to[R, l=$R_n$] (2,2)
	-- (0, 2);

\draw (2,7.5) node[xshift=0.1] {\dots} (2,5);
	\draw (2,4.5) node[xshift=0.1] {\dots} (2,4); 
\end{circuitikz}
\end{center}
per cui la tensione ai capi di una resistenza in serie sarà data da, applicando la legge di Ohm (2.1):
\[
	\begin{cases}
		i = \frac{v}{R_1 + R_2 + ...}, \quad v_1 = i R_1 \implies v_1 = v \frac{R_1}{R_1 + R_2 + ...} \\
		i = \frac{v}{R_1 + R_2 + ...}, \quad v_2 = i R_2 \implies v_2 = v \frac{R_2}{R_1 + R_2 + ...} \\
		...
	\end{cases}
\]


\par\smallskip

\noindent
\textbf{\textsf{Partitore di corrente}}

\noindent
Per le resistenze in parallelo anziché le resistenze, sommiamo le conduttanze:
$$
G_{eq} = G_1 + G_2 + ... + G_j + ... + G_n
$$

Possiamo usare questa legge per studiare il \textit{partitore di corrente}:
\begin{center}
\begin{circuitikz}
    \draw (2,-2) 
				-- (0,-2) 
				to[ european current source, i=$I(t)$] (0, 0);	

		\draw (2,0 )to[R, l=$R_1$] (2,-2);
		\draw (4,0 )to[R, l=$R_2$] (4,-2);
		
		
    \draw (4,0) node[anchor=west,xshift=0.7cm] {\dots} (6,0);
    \draw (4,-2) node[anchor=west,xshift=0.7cm] {\dots} (6,-2); 
		
		\draw (8,0 )to[R, l=$R_J$] (8,-2);
       
    \draw (8,0) node[anchor=west,xshift=0.7cm] {\dots} (10,0);
    \draw (8,-2) node[anchor=west,xshift=0.7cm] {\dots} (10,-2); 
		
		\draw (12,0 )to[R, l=$R_n$] (12,-2);

    % Draw horizontal lines at the top and bottom
    \draw (0,0) -- (4,0);
    \draw (0,-2) -- (4,-2);

    \draw (6,0) -- (8,0);
    \draw (6,-2) -- (8,-2);

    \draw (10,0) -- (12,0);
    \draw (10,-2) -- (12,-2);

\end{circuitikz}
\end{center}
per cui la corrente che scorre su ogni resistenza sarà data da, sempre applicando la legge di Ohm:
\[
	\begin{cases}
		i = v (G_1 + G_2 + ...), \quad i_1 = v G_1 \implies i_1 = i \frac{G_1}{G_1 + G_2 + ...} \\ 
		i = v (G_1 + G_2 + ...), \quad i_2 = v G_2 \implies i_2 = i \frac{G_2}{G_1 + G_2 + ...} \\ 
		...
	\end{cases}
\]

Ricordiamo quindi cosa significa effettivamente prendere la somma delle conduttanze:
$$
G_{eq} = G_1 + G_2 = \frac{1}{R_1} + \frac{1}{R_2}, \quad R_{eq} = \left( \frac{1}{R_1} + \frac{1}{R_2} \right) = \frac{R_1 R_2}{R_1 + R_2}
$$
ciò significa che il partitore di corrente si può prendere anche come:
$$
i_j = i \frac{G_j}{\sum_i^N G_i} = i \frac{\prod_{i \neq j}^N R_j}{\sum_i^N R_i} 
$$
cioè basta prendere i prodotti di tutte le resistenze non non coinvolte nel ramo su cui si calcola la corrente, al numeratore. 

\par\smallskip

\noindent
\textbf{\textsf{Principio di sovrapposizione degli effetti}}

\noindent
Vediamo questa legge, che si applica a reti \textit{lineari}:
\begin{theorem}{Principio di sovrapposizione degli effetti}
La risposta di una rete lineare che contiene più generatori indipendenti può essere ricavata considerando un singolo generatore per volta e sommando tutte le risposte così ottenute.
\end{theorem}
Facendo attenzione al fatto che parliamo di generatori \textit{indipendenti} (i generatori dipendenti fanno eccezione).

Per considerare un singolo generatore per volta dobbiamo \textit{disattivare} gli altri generatori. Ricordiamo che questo significa:
\begin{itemize}
	\item Trasformare i generatori di \textit{corrente} in \textbf{circuiti aperti};
	\item Trasformare i generatori di \textit{tensione} in \textbf{cortocircuiti}.
\end{itemize}

Notiamo che le reti che vedremo non sempre saranno lineari: in tal caso questa legge non si potrà applicare, e dovrà al limite essere applicata a versioni \textit{linearizzate} di tali reti.

\par\smallskip

\noindent
\textbf{\textsf{Teorema di Thevenin e Norton}}


\noindent
Ricordiamo questi due teoremi:
\begin{theorem}{Teorema di Thevenin}
	Un circuito lineare, rispetto ad una coppia dei suoi nodi, può essere trasformato in una combinazione di un generatore di tensione $v_T$ e una resistenza $R_{eq}$.
\end{theorem}

In particolare, la tensione $v$ è quella vista dai due terminali se si separa la rete dal resto del circuito e si mette fra loro un circuito aperto. La resistenza $R$ è semplicemente la resistenza vista fra i 2 morsetti, disattivati i generatori indipendenti. 
Notiamo che nel caso ci siano generatori dipendenti, questi vanno mantenuti e bisogna introdurre un generatore di prova (di corrente o di tensione, si calcolerà la resistenza \textit{vista} misurando l'altra grandezza e mettendola a rapporto).

\begin{theorem}{Teorema di Norton}
	Un circuito lineare, rispetto ad una coppia dei suoi nodi, può essere trasformato in una combinazione di un generatore di corrente $i_N$ e una resistenza $R_{eq}$.
\end{theorem}

In particolare, la tensione $i$ è quella vista dai due terminali se si separa la rete dal resto del circuito e si mette fra loro un cortocircuito. La resistenza $R_{eq}$ è analoga fra Thevenin e Norton.

Un corollario dei teoremi di Thevenin e Norton è che un generatore di corrente in parallelo a una resistenza equivale ad un generatore di tensione in serie alla stessa resistenza (salvo casi limite).

\par\smallskip

\noindent
\textbf{\textsf{Capacità ed induttanze}}

\noindent
Introduciamo altri due componenti circuitali, cioè le \textbf{capacità} e le \textbf{induttanze}.
\begin{itemize}
	\item Definiamo la \textit{capacità} come l'elemento passivo descritto dalla legge:
		$$
			q = Cv
		$$
		dove $C$ viene detta \textbf{capacità} e si misura in $F$ (\textit{Farad}). Dal punto di vista temporale abbiamo quindi:
		$$
		i = \frac{dq}{dt} = C \frac{dv}{dt}, \quad v = \frac{1}{C} \int_{t_0}^t i \, d\tau + v(t_0) 
		$$
		In DC, quindi, la capacità si comporta come un \textit{circuito aperto}, e la \textit{tensione} non può variare istantaneamente (porta energia).

		Riguardo alle capacità in \textit{serie}, abbiamo che queste si comportano nella maniera opposta alle resistenze, cioè:
		$$
		C_{eq} = \left( \frac{1}{C_1} + \frac{1}{C_2} + ... \right)^{-1}
		$$
		mentre per le capacità in \textit{parallelo} abbiamo:
		$$
		C_{eq} = C_1 + C_2 + ...
		$$
	\item Definiamo invece l'\textit{induttanza} come l'elemento passivo descritto dalla legge:
		$$
			\phi = Ci
		$$
		dove $L$ viene detta \textbf{induttanza} e si misura in $H$ (\textit{Henry}). $\phi$ è invece il flusso magnetico. Dal punto di vista temporale abbiamo quindi:
		$$
		v = \frac{d\phi}{dt} = L \frac{di}{dt}, \quad i = \frac{1}{L} \int_{t_0}^t v \, d\tau + i(t_0)
		$$
		In DC, quindi, l'induttanza si comporta come un \textit{cortocircuito}, e la \textit{corrente} non può variare istantaneamente (porta energia).
		
		Riguardo alle induttanze in \textit{serie}, abbiamo che queste si comportano come le resistenze, cioè:
		$$
		L_{eq} = L_1 + L_2 + ...
		$$
		mentre per le induttanze in \textit{parallelo} abbiamo:
		$$
		L_{eq} = \left( \frac{1}{L_1} + \frac{1}{L_2} + ... \right)^{-1}
		$$
\end{itemize}

\subsubsection{Transitorio del condensatore}
Dopo aver dato le basi di elettrotecnica, vediamo la situazione in cui ci troviamo quando chiudiamo un circuito con una resistenza e una capacità (il celebre circuito \textbf{RC}).

\begin{center}
	\begin{tikzpicture}
		\draw (0,0) to [ voltage source,  l=$V_S$] (0,3)
			to [ resistor, l=$R$] (3,3)
			to [ capacitor, l=$C$] (3,0)
			-- (0,0);
	\end{tikzpicture}
\end{center}

Possiamo studiare questo circuito in una varietà di modi. Scegliamo il più semplice, partendo dall'impostare l'equazione di maglia:
$$
\frac{V_s}{R} = -\frac{v}{R} + C \frac{dv}{dt} \implies \frac{V_s - v}{R} = C \frac{dv}{dt} \implies \frac{dv}{v - V_s} = - \frac{dt}{RC}
$$
separando le variabili. Risolviamo quindi intengrando da entrambi i lati:
$$
\ln (v - V_S) \Big|_{V_0}^{V(t)} = -\frac{t}{RC} \Big|_0^t \implies v(t) = V_S + (V_0 - V_S)e^{-\frac{1}{\tau}}
$$
con $\tau = RC$ detto \textit{tempo caratteristico} del circuito.

Questo dà la solita figura di crescita esponenziale della grandezza $v$, che riportiamo sul grafico:
\begin{center}
\begin{tikzpicture}
    \begin{axis}[
        xlabel={$t$},
        ylabel={$v(t)$},
        domain=-10:10, % set the x range you want,
				samples=100,
        grid=major, % add a grid
				ytick={1, 2},
				yticklabels={$V_0$, $V_S$},
				xtick={0},
				xticklabels={$0$},
				width=15cm,
				height=7cm
    ]
    \addplot[
        blue,
        thick,
				domain=-10:0
    ] {1};

    \addplot[
        blue,
        thick,
				domain=0:10
    ] {2-exp(-x)};

		\addplot[
        red,
        dotted,
				domain=0:10
    ] {2}; 
    \end{axis}
\end{tikzpicture}
\end{center}
dove le varie quantità sono:
\begin{itemize}
	\item $V_0$: tensione \textit{iniziale} presente sul condensatore;
	\item $V_S$: tensione \textit{finale} presente sul condensatore, cioè quella espressa dal generatore;
	\item $\tau$: \textit{tempo caratteristico} del circuito, determinato da $RC$.
\end{itemize}

\end{document}
